\section{Integraci\'on compleja}
\label{sec:integracion_compleja}

La integral \eqref{eq:integral_ortogonalidad} se separa en parte real e
imaginaria,
\begin{equation}
  \label{eq:separacion_integral}
  \small
  \int_{-\pi}^{\pi}e^{i(k-k')x}\dd x
    =\int_{-\pi}^{\pi}\cos((k-k')x)\dd x
    +i\int_{-\pi}^{\pi}\sin((k-k')x)\dd x,
\end{equation}
lo que da la pauta para definir la integral en el plano complejo.

\begin{definicion}[Integral compleja]
  \label{def:integral_compleja}
  Sea $f(z)$ bien definida en $D\subset\C$ y $\ga:\R\rightarrow\C$ una curva
  parametrizada contenida en $D$. Definimos
  \begin{equation}
    \label{eq:integral_compleja}
    \int_{\ga}f(z)\dd z=\int_{a}^{b}f(z(t))\dv{z}{t}\dd t .
  \end{equation}
\end{definicion}

La integral compleja es entonces una integral de l\'inea en el plano. Si
$z=x+iy$,
\begin{equation}
  \label{eq:dz}
  \dd z=\dv{z}{t}\dd t=\dv{}{t}(x+iy)\dd t
    =(\dot{x}+i\dot{y})\dd t=\dd x+i\dd y,
\end{equation}
de modo que la forma m\'as general posible es
\begin{equation}
  \label{eq:forma_general_integral}
  \int_{\ga}\qty{p\dd x+q\dd y},
  \qquad p,q\in\C .
\end{equation}

\subsection{Acotamiento}
\label{subsec:acotamiento_integral}

Se cumple
\begin{equation}
  \label{eq:cota_integral}
  \abs{\int_{\ga}f(z)\dd z}
    \le\int_{\ga}\abs{f(z)}\abs{\dv{z}{t}}\abs{\dd t} .
\end{equation}
Notando que si $z(t)=x(t)+iy(t)$,
\begin{equation}
  \label{eq:diferencial_arco_complejo}
  \abs{\dv{z}{t}}
    =\qty{\qty{\dv{x}{t}}^2+\qty{\dv{y}{t}}^2}^{1/2},
\end{equation}
que es el diferencial de longitud de arco, resulta en particular
\begin{equation}
  \label{eq:cota_longitud}
  \abs{\int_{\ga}\dd z}\le\int_{\ga}\abs{\dd z}
    =\text{longitud de arco de }\ga .
\end{equation}
Si $\ga$ es una circunferencia de radio $\rho$, entonces
\begin{equation}
  \label{eq:cota_circulo}
  \abs{\int_{\ga}\dd z}\le 2\pi\rho .
\end{equation}

\subsection{Independencia de la trayectoria}
\label{subsec:independencia_trayectoria}

Retomando \eqref{eq:forma_general_integral} con $p(z)$ y $q(z)$ complejas,
supongamos que existe $U(z)$ tal que
\begin{equation}
  \label{eq:potencial_U}
  \pdv{U}{x}=p,
  \qquad
  \pdv{U}{y}=q .
\end{equation}
Entonces, por la regla de la cadena,
\begin{equation}
  \label{eq:integral_exacta}
  \begin{aligned}
    \int_{\ga}\qty{p\dd x+q\dd y}
      &=\int_{a}^{b}\qty{\pdv{U}{x}\dv{x}{t}
        +\pdv{U}{y}\dv{y}{t}}\dd t\\
      &=\int_{a}^{b}\dv{}{t}U\dd t=U(b)-U(a).
  \end{aligned}
\end{equation}

\begin{teorema}
  \label{teo:independencia}
  La integral de l\'inea $\int_{\ga}\qty{p\dd x+q\dd y}$ definida sobre un
  dominio $D$ depende s\'olo de los puntos extremos de $\ga$ si y s\'olo si
  existe una funci\'on $U(z)=U(x,y)$ en $D$ con derivadas parciales
  \eqref{eq:potencial_U}.
\end{teorema}

\subsubsection{Consecuencias}
\label{subsubsec:consecuencias_independencia}

Tenemos una diferencial exacta si
\begin{equation}
  \label{eq:diferencial_exacta}
  f(z)\dd z=f(z)\dd x+if(z)\dd y
\end{equation}
proviene de una $F(z)$ en $D$ con
\begin{equation}
  \label{eq:F_parciales}
  \pdv{F}{x}=f(z),
  \qquad
  -i\pdv{F}{y}=f(z),
\end{equation}
es decir, con $\pdv{F}{x}=-i\pdv{F}{y}$, que por
\eqref{eq:condicion_derivada} es justamente la condici\'on de
Cauchy--Riemann. En otras palabras, $F(z)$ debe ser anal\'itica en $D$.

\begin{observacion}[Consecuencia fuerte]
  La integral de una funci\'on continua depende \'unicamente de los extremos
  de $\ga$ si y s\'olo si $f$ es la derivada de una funci\'on anal\'itica en
  $D$. En particular, si $\ga$ es cerrada,
  \begin{equation}
    \label{eq:integral_cerrada_nula}
    \oint_{\ga}f(z)\dd z=0
  \end{equation}
  siempre que $f(z)$ sea anal\'itica en el dominio encerrado por $\ga$. \'Este
  es el \textbf{teorema de Cauchy}.
\end{observacion}

\begin{ejemplo}
  \label{ej:potencias_cerrada}
  \begin{equation}
    \label{eq:potencias_cerrada}
    \oint_{\ga}(z-a)^n\dd z=0
    \qquad\forall n\ge 0,
  \end{equation}
  ya que la primitiva $F(z)=\dfrac{(z-a)^{n+1}}{n+1}$ es anal\'itica.
\end{ejemplo}

\begin{ejemplo}[El caso $n=-1$]
  \label{ej:integral_fundamental}
  Sea $\ga$ la circunferencia de radio $\rho$ centrada en $a$, parametrizada
  por
  \begin{equation}
    \label{eq:parametrizacion_circulo}
    z(\vph)=a+\rho e^{i\vph}
      =a+\rho\qty{\cos\vph+i\sin\vph},
    \quad\vph\in[0,2\pi],
  \end{equation}
  de modo que $\dd z=i\rho e^{i\vph}\dd\vph$. Entonces
  \begin{equation}
    \label{eq:integral_fundamental}
    \oint_{\abs{z-a}=\rho}\frac{\dd z}{z-a}
      =\int_{0}^{2\pi}\frac{i\rho e^{i\vph}}{\rho e^{i\vph}}\dd\vph
      =2\pi i ,
  \end{equation}
  resultado independiente de $\rho$ y central en todo lo que sigue.
\end{ejemplo}

\begin{ejemplo}
  \label{ej:integral_recta}
  Sea $\ga(t)=(1+i)t$ con $t\in[0,1]$, de modo que sobre $\ga$ se tiene
  $x=\Real\set{z}=t$ y $\dd z=(1+i)\dd t$. Entonces
  \begin{equation}
    \label{eq:integral_recta}
    \int_{\ga}x\dd z=\int_{0}^{1}t\,(1+i)\dd t
      =\eval{\frac{1}{2}(1+i)t^2}_{0}^{1}=\frac{1}{2}(1+i).
  \end{equation}
  N\'otese que el integrando $x$ no es anal\'itico, de modo que aqu\'i no
  aplica \eqref{eq:integral_cerrada_nula}.
\end{ejemplo}

\subsection{Deformaci\'on de contornos}
\label{subsec:deformacion_contornos}

Supongamos que $f(z)$ es anal\'itica en $D$ excepto en un punto $z=a$.
Consideremos dos curvas: $C_1$, deformable hasta $\ga$, y $C_2$, contenida en
el interior de $\ga$ y que encierra a $z=a$. Unimos ambas por un corte de
anchura $\ep$ para formar un contorno $C_3(\ep)$ que \emph{no} encierra a
$z=a$; entonces $f$ es anal\'itica sobre y dentro de $C_3$ y
\begin{equation}
  \label{eq:contorno_C3}
  \oint_{C_3}f(z)\dd z=0 .
\end{equation}
Al hacer $\ep\rightarrow 0$ las contribuciones del corte se cancelan y queda
\begin{equation}
  \label{eq:deformacion}
  \oint_{C_1}f(z)\dd z+\oint_{-C_2}f(z)\dd z=0
  \imp
  \oint_{\ga}f(z)\dd z=\oint_{C_2}f(z)\dd z .
\end{equation}
En particular podemos pedir que $C_2$ sea una circunferencia de radio $\rho$
centrada en la singularidad.

\begin{ejemplo}
  \label{ej:fracciones_parciales}
  \begin{equation}
    \label{eq:fracciones_parciales}
    \small
    \begin{aligned}
      \oint_{\abs{z}=2}\frac{\dd z}{z^2-1}
        &=\frac{1}{2}\oint_{\abs{z}=2}\frac{\dd z}{z-1}
         -\frac{1}{2}\oint_{\abs{z}=2}\frac{\dd z}{z+1}\\
        &=\frac{1}{2}\oint_{\abs{z-1}=\rho}\frac{\dd z}{z-1}
         -\frac{1}{2}\oint_{\abs{z+1}=\rho}\frac{\dd z}{z+1}\\
        &=\frac{1}{2}(2\pi i)-\frac{1}{2}(2\pi i)=0 ,
    \end{aligned}
  \end{equation}
  donde en el segundo paso se deform\'o cada contorno a un c\'irculo peque\~no
  alrededor de su polo y en el tercero se us\'o
  \eqref{eq:integral_fundamental}.
\end{ejemplo}

Generalizando: si $f(z)$ es anal\'itica en $D$ salvo en un conjunto finito de
puntos aislados $\set{z_j}_{j=1}^{n}$, entonces
\begin{equation}
  \label{eq:deformacion_general}
  \oint_{\ga}f(z)\dd z
    =\sum_{j=1}^{n}\oint_{B_\rho(z_j)}f(z)\dd z .
\end{equation}

\subsection{F\'ormula integral de Cauchy}
\label{subsec:formula_cauchy}

Sea $f(z)$ anal\'itica en $D$ y construyamos
\begin{equation}
  \label{eq:F_cociente}
  F(z)=\frac{f(z)-f(a)}{z-a} .
\end{equation}
Al ser $f$ anal\'itica, su desarrollo de Taylor alrededor de $z=a$ es
\begin{equation}
  \label{eq:taylor_en_a}
  f(z)=f(a)+\eval{\dv{}{z}f(z)}_{z=a}(z-a)+\dots,
\end{equation}
de modo que en \eqref{eq:F_cociente} el factor $(z-a)$ se cancela y
\begin{equation}
  \label{eq:F_analitica}
  F(z)=a_1+a_2(z-a)+\dots
\end{equation}
es anal\'itica en todo $D$, incluido $z=a$. Por el teorema de Cauchy
\eqref{eq:integral_cerrada_nula} su integral sobre un contorno cerrado se
anula:
\begin{equation}
  \label{eq:derivacion_fic}
  \begin{aligned}
    0&=\oint_{\ga}\frac{f(z)-f(a)}{z-a}\dd z\\
     &=\oint_{\ga}\frac{f(z)}{z-a}\dd z
       -f(a)\oint_{\ga}\frac{\dd z}{z-a}\\
     &=\oint_{\ga}\frac{f(z)}{z-a}\dd z-2\pi i f(a),
  \end{aligned}
\end{equation}
usando \eqref{eq:integral_fundamental} en el \'ultimo paso.

\begin{teorema}[F\'ormula integral de Cauchy]
  \label{teo:fic}
  \begin{equation}
    \label{eq:fic}
    f(a)=\frac{1}{2\pi i}\oint_{\ga}\frac{f(z)}{z-a}\dd z,
  \end{equation}
  o, escribiendo la variable de integraci\'on como $\xii$,
  \begin{equation}
    \label{eq:fic_xi}
    f(z)=\frac{1}{2\pi i}\oint_{\ga}\frac{f(\xii)}{\xii-z}\dd\xii .
  \end{equation}
\end{teorema}

Es decir, los valores de una funci\'on anal\'itica en el interior de $\ga$
quedan determinados por sus valores sobre $\ga$.

\begin{ejemplo}
  \label{ej:fic_aplicada}
  Si $g(z)=e^{2z^2}-\cos^3\qty{z^{1/2}+3z^{1/3}}$ es anal\'itica dentro de
  $\ga$ y $z=2$ est\'a en el interior, entonces
  \begin{equation}
    \label{eq:fic_aplicada}
    \oint_{\ga}\frac{g(z)}{z-2}\dd z=2\pi i\,g(2),
  \end{equation}
  sin necesidad de evaluar ninguna integral.
\end{ejemplo}

\begin{definicion}[\'Indice]
  \label{def:indice}
  \begin{equation}
    \label{eq:indice}
    I=\frac{1}{2\pi i}\oint_{\ga}\frac{\dd z}{z-a}\in\Z,
  \end{equation}
  el n\'umero de vueltas que $\ga$ da alrededor de $a$.
\end{definicion}

\subsection{F\'ormula de Cauchy generalizada}
\label{subsec:cauchy_generalizada}

Derivando \eqref{eq:fic_xi} bajo el signo integral,
\begin{equation}
  \label{eq:fic_derivada}
  \dv{}{z}f(z)=\dv{}{z}\set{\frac{1}{2\pi i}
    \oint_{\ga}\frac{f(\xii)}{\xii-z}\dd\xii}
    =\frac{1}{2\pi i}\oint_{\ga}\frac{f(\xii)}{(\xii-z)^2}\dd\xii,
\end{equation}
y en general
\begin{equation}
  \label{eq:fic_general}
  \dnv{}{z}{n}f(z)
    =\frac{n!}{2\pi i}\oint_{\ga}
      \frac{f(\xii)}{(\xii-z)^{n+1}}\dd\xii .
\end{equation}

\subsection{Desigualdad de Cauchy y teorema de Liouville}
\label{subsec:liouville}

Combinando Taylor \eqref{eq:taylor_complejo} con \eqref{eq:fic_general}, los
coeficientes del desarrollo $f(z)=\sum_n C_n(z-a)^n$ son
\begin{equation}
  \label{eq:coeficientes_fic}
  C_n=\frac{1}{2\pi i}\oint_{\ga}
    \frac{f(\xii)}{(\xii-a)^{n+1}}\dd\xii .
\end{equation}
Acotando con \eqref{eq:cota_integral} y deformando $\ga$ a la circunferencia
$\xii=a+\rho e^{i\vph}$, para la cual $\abs{\dd\xii}=\rho\dd\vph$,
\begin{equation}
  \label{eq:cota_coeficientes}
  \begin{aligned}
    \abs{C_n}
      &\le\frac{1}{2\pi}\oint_{\ga}
        \frac{\abs{f(\xii)}\abs{\dd\xii}}{\abs{\xii-a}^{n+1}}
       =\frac{1}{2\pi}\int_{0}^{2\pi}
        \frac{\abs{f(\vph)}\rho\dd\vph}{\rho^{n+1}}\\
      &=\frac{1}{2\pi}\int_{0}^{2\pi}
        \frac{\abs{f(\vph)}}{\rho^{n}}\dd\vph .
  \end{aligned}
\end{equation}
Si $f$ est\'a acotada, $\abs{f}\le M$, se obtiene la
\textbf{desigualdad de Cauchy}
\begin{equation}
  \label{eq:desigualdad_cauchy}
  \abs{C_n}\le\frac{M}{\rho^{n}} .
\end{equation}
Como $0\le\rho<R$ con $R$ el radio de convergencia, $1/R\le 1/\rho$ y el
supremo se alcanza en
\begin{equation}
  \label{eq:supremo_coeficientes}
  \abs{C_n}\le\frac{M}{R^{n}} .
\end{equation}

\begin{teorema}[Liouville]
  \label{teo:liouville}
  Toda funci\'on anal\'itica y acotada en todo el plano debe ser constante.
\end{teorema}

\begin{demostracion}
  Si $f$ es entera y acotada, $R$ puede tomarse arbitrariamente grande en
  \eqref{eq:supremo_coeficientes}, de modo que $C_n=0$ para todo $n\ge 1$ y
  s\'olo sobrevive $C_0$.
\end{demostracion}

\subsection{Singularidades removibles}
\label{subsec:singularidades_removibles}

\begin{teorema}
  \label{teo:singularidad_removible}
  Supongamos $f(z)$ anal\'itica en $D$ excepto en un punto aislado $z=a$. Una
  condici\'on necesaria y suficiente para que exista una extensi\'on
  anal\'itica de $f$ a todo $D$, incluido $z=a$, es
  \begin{equation}
    \label{eq:singularidad_removible}
    \lim_{z\rightarrow a}(z-a)f(z)=0 .
  \end{equation}
  Adem\'as, la funci\'on extendida es \'unica.
\end{teorema}

\begin{ejemplo}
  \label{ej:sinc}
  Sea $f(z)=\dfrac{\sin z}{z}$. Como
  $\lim_{z\rightarrow 0}\dfrac{\sin z}{z}=1$, se tiene
  \begin{equation}
    \label{eq:sinc}
    \lim_{z\rightarrow 0}zf(z)
      =\lim_{z\rightarrow 0}\sin z=0,
  \end{equation}
  y la singularidad en $z=0$ es removible.
\end{ejemplo}

\subsubsection{Construcci\'on iterativa de Taylor}
\label{subsubsec:construccion_taylor}

La idea anterior se puede iterar. Partiendo de
\begin{equation}
  \label{eq:F0}
  F_0(z)=F(z)=\frac{f(z)-f(a)}{z-a},
\end{equation}
que cumple $\lim_{z\rightarrow a}F(z)=\eval{\dv{}{z}f(z)}_{z=a}$ y por lo
tanto $\lim_{z\rightarrow a}(z-a)F(z)=0$, definimos recursivamente
\begin{equation}
  \label{eq:Fn_recursion}
  F_n(z)=\frac{F_{n-1}(z)-F_{n-1}(a)}{z-a},
\end{equation}
de modo que
\begin{equation}
  \label{eq:limite_Fn}
  \lim_{z\rightarrow a}F_1(z)
    =\eval{\dv{}{z}F(z)}_{z=a}
    =\frac{1}{2}\eval{\ddv{}{z}f(z)}_{z=a},
\end{equation}
y en general
\begin{equation}
  \label{eq:Fn_derivada}
  F_n(a)=\frac{1}{n!}\eval{\dnv{}{z}{n}f(z)}_{z=a} .
\end{equation}
Despejando hacia atr\'as la recursi\'on \eqref{eq:Fn_recursion} se reconstruye
el desarrollo de Taylor con residuo.

\begin{teorema}
  \label{teo:taylor_residuo}
  Si $f(z)$ es anal\'itica en una regi\'on $D$ que contiene a $z=a$, entonces
  \begin{equation}
    \label{eq:taylor_residuo}
    \small
    \begin{aligned}
      f(z)&=f(a)+\frac{1}{1!}f'(a)(z-a)
        +\frac{1}{2!}f''(a)(z-a)^2+\dots\\
        &\quad+\frac{f^{(n)}(a)}{n!}(z-a)^n
        +F_{n+1}(z)(z-a)^{n+1},
    \end{aligned}
  \end{equation}
  con $F_{n+1}(z)$ anal\'itica en $D$.
\end{teorema}

\subsection{Ceros}
\label{subsec:ceros}

Supongamos que $f(z)$ es anal\'itica en $D$ y que $f(a)=0$. Si adem\'as
\begin{equation}
  \label{eq:cero_orden_n}
  \eval{\dnv{}{z}{\ka}f(z)}_{z=a}=0
  \quad\forall\,0\le\ka\le n-1,
  \qquad
  \eval{\dnv{}{z}{n}f(z)}_{z=a}\neq 0,
\end{equation}
decimos que $z=a$ es un \textbf{cero de orden $n$}. En tal caso el desarrollo
de Taylor empieza en la potencia $n$ y podemos factorizar
\begin{equation}
  \label{eq:factorizacion_cero}
  f(z)=(z-a)^n\tilde{f}(z),
  \qquad \tilde{f}(a)\neq 0,
\end{equation}
con $\tilde{f}$ anal\'itica. Si hay varios ceros $\set{z_j}$ en $D$,
\begin{equation}
  \label{eq:factorizacion_multiple}
  f(z)=(z-z_1)^{m_1}(z-z_2)^{m_2}\dots(z-z_r)^{m_r}\tilde{f}(z),
\end{equation}
donde $m_k$ es la \textbf{multiplicidad} de cada cero y
$\sum_{k=1}^{r}m_k=m$ es el orden total.

\subsection{Polos y series de Laurent}
\label{subsec:polos_laurent}

Sea $f(z)$ anal\'itica en $D$ excepto en un conjunto de puntos aislados. En
cada uno de ellos se distinguen dos casos:
\begin{itemize}
  \item si $\lim_{z\rightarrow a}(z-a)f(z)=0$, la singularidad es
        \textbf{removible};
  \item si $\lim_{z\rightarrow a}f(z)=\infty$, tenemos un \textbf{polo}.
\end{itemize}

Supongamos que $a$ es un polo y tomemos un disco de radio $\tilde{R}<R$ donde
$f(z)\neq 0$ para $0<\abs{z-a}<\tilde{R}$. Definimos
\begin{equation}
  \label{eq:g_reciproca}
  g(z)=\frac{1}{f(z)},
\end{equation}
anal\'itica salvo en $z=a$ y con $\lim_{z\rightarrow a}g(z)=0$; por el
Teorema~\ref{teo:singularidad_removible} se extiende anal\'iticamente al
disco. Entonces $z=a$ es un cero de orden $n$ de $g$, y por
\eqref{eq:factorizacion_cero}
\begin{equation}
  \label{eq:g_factorizada}
  g(z)=(z-a)^n\tilde{g}(z),
  \qquad \tilde{g}(a)\neq 0 .
\end{equation}
En consecuencia
\begin{equation}
  \label{eq:f_polo}
  f(z)=\frac{1}{(z-a)^n}\tilde{f}(z),
  \qquad
  \tilde{f}(z)=\frac{1}{\tilde{g}(z)} .
\end{equation}
Como $\tilde{f}$ es anal\'itica, admite desarrollo de Taylor, y sustituyendo
\begin{equation}
  \label{eq:laurent_desarrollo}
  \begin{aligned}
    f(z)&=\frac{1}{(z-a)^n}\sum_{k=0}^{\infty}C_k(z-a)^k\\
        &=\frac{B_n}{(z-a)^n}+\dots+\frac{B_1}{(z-a)}
          +B_0+B_1'(z-a)+\dots
  \end{aligned}
\end{equation}

\begin{definicion}[Serie de Laurent]
  \label{def:laurent}
  En general, alrededor de una singularidad de $f(z)$ escribimos
  \begin{equation}
    \label{eq:laurent}
    f(z)=\sum_{k=-\infty}^{\infty}C_k(z-a)^k,
  \end{equation}
  que se separa en
  \begin{equation}
    \label{eq:laurent_partes}
    f(z)=\underbrace{\sum_{k=-\infty}^{-1}C_k(z-a)^k}
           _{\text{parte principal o singular}}
      +\underbrace{\sum_{k=0}^{\infty}C_k(z-a)^k}
           _{\text{parte anal\'itica}} .
  \end{equation}
\end{definicion}

Seg\'un cu\'antos t\'erminos tenga la parte principal:
\begin{itemize}
  \item \textbf{polo de orden $n$} si tiene exactamente $n$ t\'erminos;
  \item \textbf{singularidad esencial} si tiene infinitos t\'erminos.
\end{itemize}
